\section{Related work}\label{sec:related}

\paragraph{Finite element software.}
General-purpose open-source finite element libraries with Python interfaces include FEniCS/DOLFINx \cite{logg2012fenics,baratta2023dolfinx},
MFEM \cite{anderson2021mfem}, Kratos Multiphysics \cite{dadvand2010kratos}, scikit-fem \cite{gustafsson2020skfem} and SfePy \cite{cimrman2019sfepy}.
OpenSeesPy \cite{zhu2018openseespy} exposes the OpenSees framework for structural and earthquake engineering from Python.
The first group is organised around weak forms and generic discretisation, the last around ready-made structural elements.
Recently, finite element tooling built natively on PyTorch has also appeared: the TensorMesh library \cite{tensormesh2026} and the TensorGalerkin assembly algorithm of Wen et al.\ \cite{wen2026tensorgalerkin}.
\fea{} follows the structural-element tradition: elements are classes with explicit $\bm B$-matrices and stiffness, mass and internal-force
methods rather than forms assembled from a variational language.

\paragraph{Model order reduction software.}
pyMOR \cite{milk2016pymor} is a mature Python library for reduced-basis and related methods, designed around abstract operator
interfaces and parameterised models. \rom{} takes a smaller route: methods act on explicit matrices and vectors, and the
FE side only has to supply them. The methods themselves are established; \cref{sec:rom} gives the references.

\paragraph{Scope comparison.}
\Cref{tab:scope} lists the focus of each package in qualitative terms, taken from the cited publications. It does not rank the packages
and says nothing about performance.

\begin{table}[t]
\centering\footnotesize
\caption{Qualitative scope of related software, as described in the cited publications. The last row is this work.}
\label{tab:scope}
\begin{tabular}{>{\raggedright\arraybackslash}p{2.9cm}>{\raggedright\arraybackslash}p{3.4cm}>{\raggedright\arraybackslash}p{3.9cm}>{\raggedright\arraybackslash}p{3.6cm}}
\toprule
Package & Core idea & Typical use & Reduction methods \\
\midrule
FEniCS/DOLFINx \cite{logg2012fenics,baratta2023dolfinx} & automated solution of PDEs from weak forms & general PDE problems & not part of the core \\
MFEM \cite{anderson2021mfem} & scalable high-order finite elements, C++ with Python bindings & large-scale parallel simulation & not part of the core \\
Kratos \cite{dadvand2010kratos} & multi-physics framework & coupled multi-physics & separate application modules, including reduced-order work (not assessed here) \\
scikit-fem \cite{gustafsson2020skfem} & assembly of weak forms in NumPy/SciPy & research prototypes & not part of the core \\
SfePy \cite{cimrman2019sfepy} & problem description files, weak forms & general PDE problems & homogenisation, not a MOR suite \\
OpenSeesPy \cite{zhu2018openseespy} & structural-engineering framework with Python interface & structural and earthquake analysis & not part of the core \\
pyMOR \cite{milk2016pymor} & abstract operator interface for reduction & parametric reduced-basis models & reduced-basis, balanced truncation and others \\
\midrule
\fea{} + \rom{} & structural elements and ROMs on explicit arrays & structural FE with reduction in one place & \cref{sec:rom} \\
\bottomrule
\end{tabular}
\end{table}
